%=====================================================================
\chapter{Indexes, and Where They Stop Helping}\label{ch:indexes}
%=====================================================================
\locator{6}{Part VI --- Making It Work, and What Comes Next}

\begin{outcomes}
By the end of this chapter you will be able to:
\begin{tight}
  \item \textbf{Explain} what an index is, using the B-tree, without
        assuming any course in data structures.
  \item \textbf{Predict} whether a query can use a given index, including a
        composite one, from the order of its columns.
  \item \textbf{Find} the share of a table above which an index is slower than
        reading everything, by measurement.
  \item \textbf{Explain} why SQLite's optimiser chose a twenty-times slower
        plan, and \textbf{override} it.
  \item \textbf{State} what an index costs, and \textbf{decide} whether a
        proposed one is worth it.
\end{tight}
\end{outcomes}

\begin{prereq}
\chref{whodoeswhat}: the optimiser and the buffer manager. \chref{firstdatabase}:
the database is a file of 4{,}096-byte pages. \chref{joins},
\chref{subqueries} and \chref{grouping}, each of which ended with an index
making something faster --- this chapter is about why, and about where it
stops. \textbf{No data-structures course is assumed.} You need
\texttt{code/m21\_indexes.py}.
\end{prereq}

\begin{keyterms}
index $\bullet$ B-tree $\bullet$ node $\bullet$ fan-out $\bullet$ depth
$\bullet$ leaf $\bullet$ rowid $\bullet$ covering index $\bullet$ composite
index $\bullet$ leftmost prefix $\bullet$ selectivity $\bullet$ crossover
$\bullet$ sequential read $\bullet$ random read $\bullet$ \sqlkw{ANALYZE}
$\bullet$ \texttt{sqlite\_stat1} $\bullet$ STAT4 $\bullet$ \sqlkw{INDEXED BY}
$\bullet$ \sqlkw{NOT INDEXED}
\end{keyterms}

\section{What an Index Is}

A table in SQLite is stored in order of its rowid. To find the student with
roll number \texttt{2024-IT-00317} in a table with no index on roll number,
the only method is to look at every row --- a \sqlkw{SCAN}. Four chapters have
shown plans that said \sqlkw{SEARCH} instead, and every one of them was using
the same idea.

\begin{definitionbox}[Index]
An \term{index} on a column is a separate structure, kept in the same file,
that holds every value of that column \textbf{in sorted order}, each paired
with the rowid of the row it came from.

\smallskip
Because it is sorted, a value can be found without looking at most of them ---
exactly as a word is found in a dictionary without reading the dictionary.
Because each entry carries a rowid, the full row can then be fetched directly.

\smallskip
SQLite keeps every index as a \term{B-tree}. You do not need to know how one is
built to use it; you need to know one number about it, below.
\end{definitionbox}

\begin{conceptbox}[The B-tree, in one paragraph and one number]
A B-tree is a tree of pages. The top page holds a few hundred keys that divide
the whole range of values into a few hundred parts, each with a pointer to a
page below; that page divides its part into a few hundred more, and so on,
until the bottom pages --- the \term{leaves} --- hold the actual entries.

\smallskip
The number that matters is the \term{fan-out}: how many children each page
points to. A 4{,}096-byte page holding short integer keys has a fan-out of a
few hundred. So:

\begin{tabular}{@{}L{3.6cm}L{10.2cm}@{}}
1 level   & a few hundred entries \\
2 levels  & tens of thousands \\
3 levels  & millions \\
4 levels  & hundreds of millions \\
\end{tabular}

\smallskip
\textbf{Finding any value means reading one page per level} --- three or four
reads, whether the table has a thousand rows or a million. That is the whole
reason part~1 of the measurement below shows a lookup costing the same at every
size.
\end{conceptbox}

\dsfig{%
\begin{tikzpicture}[
  pg/.style={rounded corners=2pt, draw=#1, line width=0.9pt, fill=#1!8,
             text=#1!60!black, font=\scriptsize\ttfamily, minimum height=0.55cm,
             inner sep=3pt},
  lf/.style={pg=greenhl, minimum width=1.9cm},
  ln/.style={draw=neutral!60, line width=0.7pt},
  hot/.style={draw=accent, line width=1.4pt, -{Stealth[length=2.2mm]}},
  lv/.style={font=\scriptsize\itshape, text=neutral, anchor=east},
  nt/.style={font=\scriptsize\itshape, text=neutral, align=left}
]
  \node[pg=primary, minimum width=4.6cm] (r) at (0,2.4)
       {\textbar\ 250k \textbar\ 500k \textbar\ 750k \textbar};
  \node[pg=secondary, minimum width=3.0cm] (i1) at (-4.5,1.1)
       {\textbar\ 60k \textbar\ 120k \textbar\ \dots};
  \node[pg=secondary, minimum width=3.0cm] (i2) at (-0.9,1.1)
       {\textbar\ 310k \textbar\ 370k \textbar\ \dots};
  \node[font=\scriptsize, text=neutral] at (2.1,1.1) {\dots};
  \node[pg=secondary, minimum width=3.0cm] (i4) at (4.6,1.1)
       {\textbar\ 810k \textbar\ 870k \textbar\ \dots};
  \foreach \c in {i1,i2,i4}{\draw[ln] (r.south) -- (\c.north);}

  \node[lf] (l1) at (-2.6,-0.25) {317 512 $\rightarrow$ rowid};
  \node[lf] (l2) at (-0.4,-0.25) {317 513 $\rightarrow$ rowid};
  \node[font=\scriptsize, text=neutral] at (1.4,-0.25) {\dots};
  \draw[ln] (i2.south) -- (l1.north);
  \draw[ln] (i2.south) -- (l2.north);

  \node[pg=neutral, minimum width=3.6cm] (row) at (-1.5,-1.55)
       {the row itself, in the table};

  \draw[hot] ([xshift=2mm]r.south west) ++(1.3,0) -- (i2.north);
  \draw[hot] (i2.south) -- (l1.north);
  \draw[hot] (l1.south) -- (row.north);

  \node[lv] at (-6.1,2.4) {root};
  \node[lv] at (-6.1,1.1) {interior};
  \node[lv] at (-6.1,-0.25) {leaves};
  \node[lv] at (-6.1,-1.55) {table};
  \node[nt, text width=4.0cm, anchor=west] at (2.4,-0.95)
       {\textcolor{accent}{\textbf{finding k = 317\,512}}: one page at each
        level, then one jump to the row --- four reads among a million rows};
\end{tikzpicture}}%
{An index on \textit{k} over a million rows, three levels deep. Each page
divides its range among a few hundred children, so the path from the root to
any key is three pages long, and one more read fetches the row. The scan it
replaces reads every page of the table.}{btree}

\section{Measured: Where an Index Wins, and Where It Loses}

\begin{measured}[A 3{,}501$\times$ win, a 21$\times$ loss, and an optimiser that could not tell the difference]
\textsf{\footnotesize \texttt{code/m21\_indexes.py}, a table of up to
1{,}000{,}000 rows of about 80 bytes, column \textit{k} a random permutation
(so index order and table order are unrelated). SQLite 3.51.2, Intel i5-6600,
16\,GB, Windows 11, mechanical disk. Best of three.}

\medskip
\textbf{1. A point lookup} --- \texttt{SELECT v FROM t WHERE k = ?}

\begin{tabular}{@{}R{2.6cm}R{2.6cm}R{2.6cm}R{2.6cm}@{}}
\textbf{rows} & \textbf{scan} & \textbf{index} & \textbf{speed-up} \\
10{,}000    & 1.37\,ms   & 0.053\,ms & 26$\times$ \\
100{,}000   & 19.03\,ms  & 0.053\,ms & 360$\times$ \\
1{,}000{,}000 & 199.90\,ms & 0.057\,ms & \textbf{3{,}501$\times$} \\
\end{tabular}

\smallskip
The scan grows in proportion to the table. The index barely moves: a hundred
times the rows added one level, at most, to a three-level tree.

\medskip
\textbf{2. The crossover} --- \texttt{SUM(length(v)) WHERE k BETWEEN 0 AND
}$x$, each plan forced, over a growing share of a million rows

\begin{tabular}{@{}R{1.6cm}R{2.2cm}R{2.4cm}R{2.4cm}L{1.6cm}L{2.6cm}@{}}
\textbf{share} & \textbf{rows} & \textbf{via index} & \textbf{full scan}
& \textbf{faster} & \textbf{optimiser chose} \\
0.1\% & 1{,}000    & 7.1\,ms    & 178.1\,ms & index & index \\
1\%   & 10{,}000   & 90.6\,ms   & 207.0\,ms & index & index \\
2\%   & 20{,}000   & 182.0\,ms  & 210.0\,ms & index & index \\
5\%   & 50{,}000   & 436.4\,ms  & 190.3\,ms & scan  & \textbf{index} \\
10\%  & 100{,}000  & 784.3\,ms  & 202.5\,ms & scan  & \textbf{index} \\
20\%  & 200{,}000  & 1{,}549.9\,ms & 220.1\,ms & scan & \textbf{index} \\
50\%  & 500{,}000  & 3{,}955.3\,ms & 287.5\,ms & scan & \textbf{index} \\
100\% & 1{,}000{,}000 & 7{,}934.8\,ms & 381.0\,ms & scan & \textbf{index} \\
\end{tabular}

\smallskip
\textbf{The index stops winning at about 2\% of the table.} Below that it reads
only the rows it needs. Above it, every matching row is a separate jump to a
different page of the table --- and a scan, which reads each page once and in
order, wins by more as more rows are wanted. At the whole table the index plan
was \textbf{21$\times$ slower}.

\medskip
\textbf{3. Why the optimiser chose the index every time}

\begin{tabular}{@{}L{7.8cm}L{6.0cm}@{}}
STAT4 histograms compiled into this SQLite & \textbf{no} \\
what \sqlkw{ANALYZE} recorded in \texttt{sqlite\_stat1} & \texttt{'1000000 1'} \\
the whole table, as written & 7{,}952.4\,ms (\texttt{SEARCH \dots\ ik}) \\
the whole table, with \texttt{+k} & 407.2\,ms (\texttt{SCAN t}) \\
\end{tabular}

\smallskip
\texttt{'1000000 1'} is everything the optimiser knew: a million rows, one row
per value of \textit{k}. It knew \textit{k} was unique. It did \emph{not} know
how the values were spread, so it could not tell \texttt{BETWEEN 0 AND 999}
from \texttt{BETWEEN 0 AND 999999}, and treated every range on an indexed
column as selective. The histograms that would have told it --- STAT4 --- are
not compiled into the SQLite that Python ships.

\medskip
\textbf{4. What indexes cost} --- 100{,}000 rows; insert 10{,}000 more; median
of 3

\begin{tabular}{@{}L{3.2cm}R{3.0cm}R{3.0cm}R{3.2cm}@{}}
& \textbf{file bytes} & \textbf{insert} & \textbf{(range)} \\
0 indexes & 8{,}232{,}960  & 71.5\,ms  & 47.6--79.0 \\
1 on \textit{k} & 9{,}383{,}936 & 64.8\,ms & 62.0--69.2 \\
+ 1 on \textit{v} (68-byte text) & 17{,}272{,}832 & 137.9\,ms & 123.8--150.2 \\
+ 1 on (\textit{k}, \textit{v}) & 25{,}485{,}312 & 141.5\,ms & 135.2--168.9 \\
\end{tabular}

\smallskip
\textbf{Read the ranges first.} Zero and one index overlap: the new rows arrived
in key order, so each new entry went on the right-hand end of the index --- the
cheapest case, as for roll numbers issued in sequence. The index on the wide
text column is the one that cost: it doubled the insert time and added more
bytes than the table had.
\end{measured}

\begin{alertbox}[What the crossover means in practice]
\begin{tight}
  \item \textbf{An index is for finding a few rows, and ``a few'' must be
        measured.} About 2\% for this million-row table. On the seeded
        \textit{enrolment} table --- 60{,}037 narrower rows --- an index on
        \textit{grade} selecting 21\% ran in 8.94\,ms against a scan's
        9.40\,ms: a tie, so the crossover there sits near a fifth. The share
        moves with table size, row width and cache; the shape --- a share past
        which the index loses --- does not.
  \item \textbf{The optimiser can be wrong, and here it was wrong every time
        past 2\%} --- not through a bug, but because it was never told how the
        data was spread. ``The database will choose the best plan'' assumes it
        has the information to.
  \item \textbf{Two ways to overrule it.} Write \texttt{+k} instead of
        \texttt{k} in the condition, which hides the column from the index for
        that one query; or name the table \texttt{t NOT INDEXED}. Use either
        only on a query you have measured.
  \item \textbf{Low-selectivity columns rarely deserve an index.} An index on
        \textit{programme} (four values) or \textit{grade} (six) selects a
        quarter or a fifth of the table for every value --- at or past the
        crossover. Measured above, the grade index only tied a scan: it bought
        nothing, and it still costs bytes and a write on every change.
\end{tight}
\end{alertbox}

\section{Using an Index Well}

\begin{definitionbox}[Composite index, leftmost prefix, covering index]
A \term{composite index} on $(a, b)$ is sorted by $a$, and by $b$ within each
$a$ --- like a telephone directory sorted by surname, then first name.

\smallskip
It can serve a condition on $a$, or on $a$ and $b$ together, but \textbf{not
on $b$ alone}: the $b$ values are scattered through every $a$. This is the
\term{leftmost-prefix} rule, and \chref{subqueries}'s measurement turned on it
--- \textit{enrolment}'s primary key is (roll\_no, course\_code, \dots), so it
could find a student's enrolments and could not find a course's.

\smallskip
A \term{covering index} contains every column a query needs, so the table
itself is never read. \texttt{USING COVERING INDEX} in a plan means exactly
that, and it removes the per-row jump that loses past the crossover.
\end{definitionbox}

\begin{examplebox}[Ordering: what an index gives for free, and what it cannot]
From part~3 of the measurement, on 100{,}000 rows with an index on \textit{k}:

\medskip
\begin{tabular}{@{}L{6.4cm}R{2.0cm}L{5.4cm}@{}}
\texttt{ORDER BY k LIMIT 10} & 0.08\,ms & reads in index order, stops at 10 \\
\texttt{ORDER BY v LIMIT 10} & 23.87\,ms & \texttt{TEMP B-TREE}: sorts every row \\
\texttt{ORDER BY k \% 7, k LIMIT 10} & 24.74\,ms & an expression: sorts too \\
\end{tabular}

\smallskip
An index is a sorted copy of \emph{one} ordering. A query that wants that
order gets it for nothing --- and, with \sqlkw{LIMIT}, stops after ten rows
instead of reading a hundred thousand. Any other order, \textbf{including an
expression of the indexed column itself}, means sorting everything first. The
third row is the one people do not expect.
\end{examplebox}

\begin{pitfall}
\begin{tight}
  \item \textbf{Indexing every column.} Each is another B-tree kept in step
        on every write. Measured: the text index doubled insert time.
  \item \textbf{Indexing a column with few values.} Every lookup selects a
        large share --- past the crossover.
  \item \textbf{Expecting an index on $(a, b)$ to serve a condition on $b$.}
        Leftmost prefix only.
  \item \textbf{Wrapping the indexed column in an expression.}
        \texttt{WHERE lower(name) = 'ali'} cannot use an index on
        \textit{name}. (SQLite allows an index on the expression itself.)
  \item \textbf{Trusting the optimiser's choice without measuring.} It picked
        a 21$\times$ slower plan here, after \sqlkw{ANALYZE}.
  \item \textbf{Measuring an index on a table whose order matches the
        column.} Then ``random'' lookups are sequential and the crossover
        vanishes. The harness shuffles \textit{k} for exactly this reason.
\end{tight}
\end{pitfall}

\begin{lab}[Find the crossover on your own machine]
\begin{lstlisting}[language=bash]
# --- 1. The measured version first: it builds a million rows --------
python code/m21_indexes.py --quick
\end{lstlisting}
\begin{lstlisting}[language=sql]
-- --- 2. On the university database, the indexes you already have -
.open university.db
SELECT name, tbl_name, sql FROM sqlite_master WHERE type = 'index';
-- Some have no sql: they were made by PRIMARY KEY and UNIQUE.

-- --- 3. Leftmost prefix, seen in a plan --------------------------
EXPLAIN QUERY PLAN SELECT * FROM enrolment WHERE roll_no = 'x';
EXPLAIN QUERY PLAN SELECT * FROM enrolment WHERE course_code = 'x';
-- Same primary-key index. One SEARCH, one SCAN. Why?

-- --- 4. A covering index ----------------------------------------
CREATE INDEX idx_enr_cg ON enrolment(course_code, grade);
EXPLAIN QUERY PLAN
SELECT grade, COUNT(*) FROM enrolment WHERE course_code = 'CS-101'
 GROUP BY grade;
-- Find COVERING. Which table was never read?

-- --- 5. An index on a column with six values --------------------
CREATE INDEX idx_enr_grade ON enrolment(grade);
.timer on
SELECT COUNT(*), AVG(section_no) FROM enrolment WHERE grade = 'B';
SELECT COUNT(*), AVG(section_no) FROM enrolment WHERE +grade = 'B';
-- Which was faster? What share of the table did 'B' select?

-- --- 6. What ANALYZE knows --------------------------------------
ANALYZE;
SELECT * FROM sqlite_stat1;
-- Read each stat string: rows, then rows per key value.

-- --- 7. An order the index cannot give --------------------------
EXPLAIN QUERY PLAN SELECT * FROM student ORDER BY name LIMIT 10;
CREATE INDEX idx_student_name ON student(name);
EXPLAIN QUERY PLAN SELECT * FROM student ORDER BY name LIMIT 10;
EXPLAIN QUERY PLAN SELECT * FROM student ORDER BY lower(name) LIMIT 10;
\end{lstlisting}
\begin{lstlisting}[language=bash]
# --- 8. The full version, a million rows (about a minute) -----------
python code/m21_indexes.py
\end{lstlisting}

\smallskip
\textbf{What to notice.} Step 5 is the crossover in your own schema. Grade B
selects 21\% of enrolments, and on the seeded table the two forms come out
level --- 8.94 against 9.40\,ms on the measuring machine --- while the
optimiser chooses the index without hesitation. Compare that with the 2\% the
million-row table gave: the same index, the same optimiser, a crossover ten
times further out, because a small table is cheap to scan. \textbf{The rule is
not ``2\%''. The rule is that the share exists and has to be measured.}
\end{lab}

\begin{checkpoint}
\begin{tightnum}
  \item A lookup took 0.053\,ms at 10{,}000 rows and 0.057\,ms at a million,
        while a scan went from 1.37 to 199.90\,ms. Explain both using the
        fan-out and depth of a B-tree, without assuming any data-structures
        course.
  \item The optimiser chose the index for a query that selected the whole
        table, after \sqlkw{ANALYZE}, and it was 21$\times$ slower. Explain
        what information it lacked, where that information would have come
        from, and two ways to make it scan instead.
  \item An index on \textit{enrolment}(roll\_no, course\_code) exists. For each
        of \texttt{WHERE roll\_no = ?}, \texttt{WHERE course\_code = ?} and
        \texttt{WHERE roll\_no = ? AND course\_code = ?}, say whether it can be
        used and why.
\end{tightnum}
\end{checkpoint}

\begin{chaptersummary}
\begin{tight}
  \item An \textbf{index} is every value of a column, in sorted order, each
        with the rowid of its row --- kept by SQLite as a \textbf{B-tree}.
  \item A B-tree's \textbf{fan-out} is a few hundred, so a million entries fit
        in three levels and a lookup reads one page per level.
  \item \textbf{Measured:} a point lookup was \textbf{3{,}501$\times$} faster
        with an index at a million rows, and cost the same at every size.
  \item \textbf{Measured: the index stopped winning at about 2\%} of the table.
        Past that, a jump per row loses to reading every page once, in order;
        at the whole table the index was \textbf{21$\times$ slower}.
  \item \textbf{Measured: the optimiser chose the index at every share}, after
        \sqlkw{ANALYZE}, because \texttt{sqlite\_stat1} records only row counts
        and STAT4 histograms are not compiled into Python's SQLite.
  \item Override a bad plan with \texttt{+column} or \sqlkw{NOT INDEXED} ---
        on a query you have measured.
  \item A \textbf{composite} index serves only its \textbf{leftmost prefix}; a
        \textbf{covering} index means the table is never read.
  \item An index gives its own order free, including for \sqlkw{LIMIT}; any
        other order, even an expression of the indexed column, costs a sort.
  \item \textbf{Measured:} indexes cost bytes and writes; a wide text index
        doubled insert time, while an integer index on in-order keys cost
        nothing measurable.
\end{tight}
\end{chaptersummary}

\begin{reviewq}
  \item \textbf{(MCQ)} An index on $(a, b)$ can be used for a condition on:
        (a)~$b$ alone; (b)~$a$ alone; (c)~either; (d)~neither.
  \item \textbf{(MCQ)} \texttt{USING COVERING INDEX} in a plan means:
        (a)~the index is unique; (b)~the table was not read; (c)~the index
        covers the whole table; (d)~a temporary index was built.
  \item \textbf{(Short)} Why does a query that wants half a table run faster as
        a scan than through an index?
  \item \textbf{(Short)} What does \texttt{sqlite\_stat1}'s string
        \texttt{'1000000 1'} tell the optimiser, and what does it not?
  \item \textbf{(Short)} Why is an index on \textit{grade} unlikely to help?
  \item \textbf{(Applied)} For the university schema, propose the three indexes
        you would add, justify each with a query from Part~IV that it serves
        and the plan it would change, and name one index you would refuse to
        add, with the reason.
  \item \textbf{(Applied)} Re-run part~2 of the harness with \textit{k} in the
        same order as the rowid instead of shuffled (change one line), and
        report what happens to the crossover. Explain the result in terms of
        sequential and random reads, and say why the shuffled version is the
        honest measurement.
\end{reviewq}
